\SourceLeaf{20}
\par
Ceterum adhuc observamus, quum substitutis pro $x,y,z$ valoribus per $p,q$ expressis, functio $W$ necessario fieri debeat identice $=0$, etiam identice i. e. independenter a valoribus ipsarum $\mathrm{d}p,\mathrm{d}q$ fieri debere
\[
0=(\lambda T+\mu U+\nu V)\mathrm{d}p+(\lambda'T+\mu'U+\nu'V)\mathrm{d}q
\]
sive haberi
\[
\begin{aligned}
\lambda T+\mu U+\nu V&=0\\
\lambda'T+\mu'U+\nu'V&=0
\end{aligned}
\]
Hinc sequitur, quantitates $\mu\nu'-\nu\mu'$, $\nu\lambda'-\lambda\nu'$, $\lambda\mu'-\mu\lambda'$ resp. ipsis $T,U,V$, sive cosinibus angulorum $QX,QY,QZ$ proportionales evadere, quod iam e supra dictis, sed remanente signorum ambiguitate, colligere licuerat.

\Article{11}
\par
Ab his disquisitionibus generalibus ad corpora sphaeroidica elliptica descendimus, quorum caussa illae fuerant susceptae. Initio abscissarum in corporis centro sumto, semiaxibusque per $A,B,C$ designatis, aequatio superficiei erit
\[
\frac{xx}{AA}+\frac{yy}{BB}+\frac{zz}{CC}=1
\]
\par
Statuemus itaque $W=\frac{xx}{AA}+\frac{yy}{BB}+\frac{zz}{CC}-1$, unde patet, pro omnibus punctis intra corpus $W$ obtinere valores negativos, positivos autem pro omnibus punctis extra corpus. Porro erit $T=\frac{2x}{AA}$, $U=\frac{2y}{BB}$, $V=\frac{2z}{CC}$; statuendo itaque
\[
\BareRoot{(\frac{xx}{A^4}+\frac{yy}{B^4}+\frac{zz}{C^4})}=\psi
\]
erit
\[
\cos QX=\frac{x}{\psi AA},\quad\cos QY=\frac{y}{\psi BB},\quad\cos QZ=\frac{z}{\psi CC}
\]
\[
\cos QM=\frac{1}{\psi r}(\frac{(a-x)x}{AA}+\frac{(b-y)y}{BB}+\frac{(c-z)z}{CC})
\]

\Article{12}
\par
Iam introducamus duas indeterminatas $p,q$ tales ut fiat
\[
\begin{aligned}
x&=A\cos p\\
y&=B\sin p.\cos q\\
z&=C\sin p.\sin q
\end{aligned}
\]
